%=============================================================================!
% 7.S  WHAT WE NOW HAVE                                                       !
%=============================================================================!
\unnumbered{What we now have}

The Legendre transform of the Lagrangian in the velocities is the
Hamiltonian $H(q, p) = \sum_k p_k\dot{q}_k - \Lag$, which for a kinetic
energy quadratic in the velocities is the energy $K + U$ written in terms
of the coordinates and their momenta. Hamilton's equations, $\dot{q} =
\partial H/\partial p$ and $\dot{p} = -\partial H/\partial q$, move every
state of phase space along a flow that runs on the contours of $H$; for
the pendulum the contours show the swinging and the circulating motions
and the separatrix between them, along which the bob creeps towards the
top as $\mathrm{e}^{-\omega_0t}$. Any quantity changes at the rate $\pb{A}
{H} = \Liou A$, and the
Taylor series in the Liouville operator gives the motion from the
starting state, at every time for the ball and the spring and over a
short time in general.

The flow has no divergence, so it keeps
the area of every region of phase space, and its Jacobian determinant is
1 (Liouville's theorem); drag shrinks areas as $\mathrm{e}^{-\gamma t}$.

Every symmetry of the Lagrangian under a change that can be made as small
as we like, and the passing of time, gives a conserved quantity
(Noether's theorem): the energy from
the passing of time, the momentum from moving everything, the angular
momentum from turning everything, with no assumption about pairs of
forces; in a periodic box the momentum survives and the angular momentum
does not. For a time-independent Hamiltonian that is even in the momenta,
reversing every momentum makes the system retrace its path. Drag breaks
this reversal symmetry.

A method of computing motion is a step map:
the forward Euler method inflates the areas and the energy of a spring by
$1 + \omega^2\dt^2$ at every step, while the symplectic Euler method, a kick
followed by a drift, keeps area and holds the energy within a band.

Chapter~8 asks where the potential energy $U$ of atoms comes from, and
Chapter~9 returns to step maps and builds the methods with which
molecular dynamics is done.
